\section{Comparative Selection Rules and Symmetry Breaking}
\label{sec:comparison}

\subsection{Rule of Mutual Exclusion}
Both $D_{6h}$ and $Y_h$ are centrosymmetric ($i \in G$). The electric dipole transition operator $\hat{\bm{\mu}} = e\mathbf{r}$ transforms as an odd function under inversion ($u$), while the polarizability tensor $\bm{\alpha}$ transforms as an even function ($g$). Consequently:
\begin{equation}
\Gamma(\hat{\bm{\mu}}) \subset \bigoplus \text{Irrep}_u, \quad \Gamma(\bm{\alpha}) \subset \bigoplus \text{Irrep}_g.
\end{equation}
For $Y_h$, $\Gamma_{\text{dipole}} = T_{1u}$ and $\Gamma_{\text{pol}} = A_g \oplus H_g$. Thus, no single vibrational mode in $\text{C}_{60}$ can be simultaneously IR- and Raman-active. Out of 174 normal modes, only the 4 four-fold degenerate $T_{1u}$ modes appear in first-order IR spectra, while only the $2A_g$ and $8H_g$ modes appear in Raman scattering.

\subsection{Symmetry Descent and Jahn--Teller Instability}
Degradation of icosahedral symmetry to subgroups ($Y_h \to D_{5d} \to C_{5v}$) lifts the degeneracy of the 4D ($G$) and 5D ($H$) vibrational and electronic levels. In ionic fullerenes ($\text{C}_{60}^-$), the degenerate ground state undergoes spontaneous geometric distortion via the $H_g$ vibrational coupling modes.
